\documentclass[11pt,a4paper]{article}
\usepackage[left=2.2cm,right=2.2cm,top=2.2cm,bottom=2.2cm]{geometry}
\usepackage[T1]{fontenc}
\usepackage{charter}
\usepackage{booktabs}
\usepackage{amsmath}
\usepackage{xcolor}
\usepackage{listings}
\usepackage[colorlinks=true,linkcolor=blue,urlcolor=blue]{hyperref}
\usepackage{caption}

\definecolor{rscblue}{RGB}{0,83,156}
\definecolor{codebg}{RGB}{246,248,249}
\newcommand{\logS}{\ensuremath{\log S}}

\lstset{
  basicstyle=\ttfamily\scriptsize,
  backgroundcolor=\color{codebg},
  frame=single, rulecolor=\color{black!25},
  breaklines=true, columns=fullflexible,
  keywordstyle=\color{rscblue}, commentstyle=\color{black!55}\itshape,
  language=Python, showstringspaces=false, xleftmargin=4pt, framexleftmargin=4pt,
}

\renewcommand{\thetable}{S\arabic{table}}
\renewcommand{\thesection}{S\arabic{section}}

\title{\vspace{-1.5cm}\large\textbf{Supplementary Information}\\[4pt]
\normalsize Notation matters: cross-representation inconsistency in chemistry\\
language models and its mechanistic origins}
\author{\normalsize Animesh Mishra}
\date{\small Manuscript DD-COM-05-2026-000309}

\begin{document}
\maketitle
\vspace{-0.6cm}
\noindent\rule{\textwidth}{0.5pt}

\noindent This document provides the verification material referenced in the response to referees: the
corrected structural classification and the script that generates it, the model validity screening
results, the full per-layer intervention distribution and controls, the robustness analyses, and the
outcome of every additional dataset evaluated, including those that failed screening and are therefore
not reported in the main text.

\section{Corrected structural classification}

The structural-class assignment in the submitted version used
\texttt{re.search(r'[FClBrI]', s)} for the halogen rule. In regular-expression syntax the square
brackets define a character class matching any single character among \texttt{F}, \texttt{C},
\texttt{l}, \texttt{B}, \texttt{r}, \texttt{I}; because \texttt{C} occurs in nearly all organic
SMILES, most molecules matched. The corrected rule tests for the element tokens explicitly.

Because assignment is priority-ordered, the defect can only affect classes evaluated \emph{after}
the halogen rule. Table~\ref{tab:cls} gives the full comparison over all 1128 ESOL molecules.

\begin{table}[h]
\centering
\caption{Class assignment under the original and corrected classifiers, all 1128 ESOL molecules.
The three classes matched upstream of the corrected rule have identical membership.}
\label{tab:cls}
\small
\begin{tabular}{lrrrl}
\toprule
Class & Original & Corrected & $\Delta$ & Position in priority chain \\
\midrule
heterocyclic      & 433 & 433 & $0$    & 1st --- upstream, unaffected \\
aromatic          & 293 & 293 & $0$    & 2nd --- upstream, unaffected \\
carboxylic/ester  &  51 &  51 & $0$    & 3rd --- upstream, unaffected \\
halogenated       & 351 &  76 & $-275$ & 4th --- the corrected rule \\
alcohol/phenol    &   0 & 179 & $+179$ & downstream \\
aliphatic         &   0 &  89 & $+89$  & downstream \\
amine             &   0 &   7 & $+7$   & downstream \\
\midrule
\multicolumn{4}{l}{Molecules whose class changes} & 275 of 1128 (24.4\%) \\
\bottomrule
\end{tabular}
\end{table}

\noindent The script below reproduces Table~\ref{tab:cls} and recomputes all class-stratified
statistics from the deposited Phase~1 predictions, without re-running model inference.

\begin{lstlisting}
def classify_original(smiles):          # as submitted
    s = smiles
    if re.search(r'[nNsS].*\d|\d.*[nNsS]', s):            return "heterocyclic"
    if re.search(r'[a-z]', s) and re.search(r'\d', s):    return "aromatic"
    if 'C(=O)O' in s or 'C(O)=O' in s or 'OC(=O)' in s:   return "carboxylic/ester"
    if re.search(r'[FClBrI]', s):                          return "halogenated"   # DEFECT
    if re.search(r'(?<![a-z])O(?![a-z=\(])', s):          return "alcohol/phenol"
    if re.search(r'(?<![a-z])N(?![a-z=\(+])', s):         return "amine"
    return "aliphatic"

def classify(smiles):                   # corrected
    s = smiles
    if re.search(r'[nNsS].*\d|\d.*[nNsS]', s):            return "heterocyclic"
    if re.search(r'[a-z]', s) and re.search(r'\d', s):    return "aromatic"
    if 'C(=O)O' in s or 'C(O)=O' in s or 'OC(=O)' in s:   return "carboxylic/ester"
    if re.search(r'Cl|Br|F|I', s):                         return "halogenated"   # FIXED
    if re.search(r'(?<![a-z])O(?![a-z=\(])', s):          return "alcohol/phenol"
    if re.search(r'(?<![a-z])N(?![a-z=\(+])', s):         return "amine"
    return "aliphatic"
\end{lstlisting}

\noindent Full script: \texttt{09\_reclassify\_phase1.py} in the deposited repository. Run as
\texttt{python3 09\_reclassify\_phase1.py phase1\_full\_results.json}.

\section{Model validity screening}

A model that emits a near-constant value irrespective of input will show a sign-flip rate determined
by the sign of that constant rather than by any property of the notation, and will appear spuriously
notation-robust. Every model is therefore screened on two per-notation diagnostics: the modal
fraction (share of predictions taking the single most frequent value) and the correlation between
predictions and ground truth. A low count of \emph{distinct} values is not itself evidence of
degeneracy, since a model answering on a coarse grid legitimately reuses values.

\begin{table}[h]
\centering
\caption{Validity screening. Only the model reported in the main text passes.}
\label{tab:valid}
\small
\begin{tabular}{llrrl}
\toprule
Model & Family / precision & Modal frac. & Pearson $r$ & Outcome \\
\midrule
qwen3.6:35b-a3b   & Qwen3 MoE, native & 21.1\% & $+0.755$ & \textbf{pass} \\
nous-hermes2:34b  & Llama, 34\,B      & 88.8\% & $+0.197$ & fail \\
qwen2.5:7b        & Qwen2.5, 7\,B     & 42.5\% & $+0.055$ & fail \\
hauhau-q4         & Qwen3 MoE, 4-bit  & ---    & ---      & fail (4/40 parseable) \\
hauhau-iq2        & Qwen3 MoE, 2-bit  & 30.8\% & $-0.177$ & fail \\
\bottomrule
\end{tabular}
\end{table}

\noindent \textbf{Retrospective validation of the reported model.} Applied to the deposited Phase~1
predictions, Qwen3.6 gives modal fractions of 19.9\% (SMILES), 13.0\% (IUPAC), 20.1\% (InChI) and
23.7\% (SELFIES), with Pearson $r = 0.726$, $0.774$, $0.753$, $0.617$ and Spearman
$\rho = 0.776$, $0.831$, $0.829$, $0.662$ (all $p < 10^{-100}$).

\noindent \textbf{Discarded run.} nous-hermes2:34b returned a single identical value for all 181
SELFIES prompts and the same value for 88.8\% of SMILES prompts. Read naively it yields 98.9\%
cross-notation inconsistency and a 0\% SELFIES sign-flip rate --- both artifacts of differing
constants per notation. These figures are not reported in the main text.

\section{Full per-layer intervention distribution and controls}

Table~\ref{tab:layers} reports every intermediate layer before any selection, answering the concern
that selecting the maximum-improvement layer guarantees a non-negative result.

\begin{table}[h]
\centering
\caption{Per-layer intervention effects and controls, all 1051 augmented-SMILES pairs.
Improvement in $\logS$; positive values narrow the prediction gap.}
\label{tab:layers}
\small
\begin{tabular}{lrrrrr}
\toprule
Layer & True partner & Random molecule & Reversed & Canonical & \% pairs degraded \\
\midrule
$L_1$ & $-0.006$ & $-0.076$ & $+0.008$ & $0.000000$ & 49.1\% \\
$L_2$ & $-0.012$ & $-0.223$ & $-0.086$ & $0.000000$ & 50.0\% \\
$L_3$ & $+0.089$ & $-0.458$ & $-0.092$ & $0.000000$ & 44.1\% \\
$L_4$ & $+0.038$ & $-0.688$ & $-0.009$ & $0.000000$ & 45.7\% \\
$L_5$ & $+0.096$ & $-0.826$ & $+0.000$ & $0.000000$ & 42.9\% \\
\bottomrule
\end{tabular}
\end{table}

\noindent \textbf{Off-manifold check.} The canonical control replaces a molecule's CLS state with its
own. It produces a deviation of exactly zero at every layer for all 1051 pairs (0 pairs nonzero),
excluding the hook mechanism, CLS substitution as such, and the Ridge head as sources of the observed
effects. This does not establish that cross-molecule substitution remains on the activation manifold,
and the main text states that limitation.

\noindent \textbf{Specificity.} True-partner substitution exceeds random-molecule substitution at
every layer, with the margin growing monotonically with depth: $+0.070$, $+0.211$, $+0.548$, $+0.726$,
$+0.922\,\logS$ at $L_1$--$L_5$ (Wilcoxon $p$ from $3.2\times10^{-10}$ to $1.1\times10^{-80}$;
Spearman $\rho = 1.000$ for margin against depth).

\noindent \textbf{Degenerate positions.} The embedding position is identical for every molecule
(cosine $1.0000$, s.d.\ $0.0000$), since the pre-attention CLS state does not depend on the input.
The final position is tautological as an intervention: the prediction head reads it, so replacing it
substitutes the source molecule's prediction wholesale, and the resulting ``improvement'' equals the
unpatched gap to within $5\times10^{-5}$ for all 1051 pairs. Both are excluded.

\section{Robustness of the behavioural effect}

\begin{table}[h]
\centering
\caption{Sensitivity of the inconsistency rate to the spread threshold ($n=1072$), and to
reformulating the task as classification ($n=790$ complete cases).}
\label{tab:robust}
\small
\begin{tabular}{lr@{\hskip 2cm}lr}
\toprule
Threshold ($\logS$) & \% inconsistent & Classification scheme & \% disagreeing \\
\midrule
0.10 & 94.5\% & 4-band & 71.8\% \\
0.25 & 93.6\% & 3-band & 45.9\% \\
0.50 & 88.0\% & 2-band (binary, $\logS=-2$) & 27.8\% \\
0.75 & 80.9\% & & \\
1.00 & 61.0\% & \multicolumn{2}{l}{Band accuracy, binary: 82--90\% across notations} \\
1.50 & 46.9\% & & \\
2.00 & 26.3\% & & \\
\bottomrule
\end{tabular}
\end{table}

\noindent \textbf{Metric.} The standard deviation across all four notation-specific predictions has
mean $0.830\,\logS$ (median $0.700$) and ranks molecules almost identically to the range
(Spearman $\rho = 0.987$, $p < 10^{-300}$, $n = 790$). Under the standard deviation on the matched
set, halogenated compounds are marginally more inconsistent than carboxylic/ester compounds
($1.094$ against $1.001$), so the identification of a single most-inconsistent class is sensitive to
metric and subset; the main text states the carboxylic/ester result as holding for the full sample
under the range metric.

\noindent \textbf{Non-language-model baseline.} A RandomForest on eight RDKit descriptors
(5-fold CV MAE $0.237\,\logS$) shows a mean prediction spread of $0.105\,\logS$ across four string
routes to the same molecule, against $1.686\,\logS$ for the language model, with 9.8\% of molecules
exceeding the $0.5\,\logS$ threshold against 88.0\%. For 79.2\% of molecules every string route
yields a bit-identical Morgan fingerprint and therefore an identical prediction.

\section{Additional datasets: all outcomes}

\begin{table}[h]
\centering
\caption{Every additional dataset evaluated, including those that failed the validity screen and are
therefore not reported as evidence in the main text.}
\label{tab:datasets}
\small
\begin{tabular}{llll}
\toprule
Dataset & Task & Screen & Outcome \\
\midrule
FreeSolv      & regression     & pass ($r = 0.646$, $0.599$, $0.460$) & effect replicates; see below \\
Lipophilicity & regression     & fail ($r = 0.180$, $0.163$, $0.200$) & not reported \\
BBBP          & classification & fail (below majority class)          & not reported \\
BACE          & classification & fail (below majority class)          & not reported \\
\bottomrule
\end{tabular}
\end{table}

\noindent \textbf{FreeSolv replication.} 81.9\% of molecules exceed a 0.5-unit three-notation spread,
against 83.6\% for the same three notations on ESOL. Mean spread is $3.364$\,kcal\,mol$^{-1}$ and
$1.404\,\logS$ respectively; normalised by the range of each target this is 14.3\% and 13.0\%.
Notation ranking does not transfer: InChI has the lowest error on ESOL and the highest of the three
on FreeSolv (MAE $3.054$ against $2.596$ for SMILES).

\noindent \textbf{Classification datasets.} On BACE the model answered positive for 88.7\% of SMILES
prompts but 8.0\% of SELFIES prompts, at 50.0\% accuracy against a 60.0\% majority class. Read
naively this gives 90.0\% cross-notation label disagreement. That figure is not reported, because a
model that cannot perform a task cannot be used to measure notation robustness on it. Whether the
effect extends to molecular classification therefore remains open.

\vspace{0.6cm}
\noindent\rule{\textwidth}{0.5pt}

\noindent\small All scripts, per-molecule predictions and control outputs are deposited with the
article; see the Data availability statement.

\end{document}
